\subsection{Período Mínimo con la Función de Prefijos (KMP)}

\graybold{Preguntas guía:}

\begin{itemize}
\item ¿Me piden si una cadena es la repetición exacta de un bloque más corto, o cuántas veces se repite el bloque más pequeño posible?
\item ¿Necesito el período mínimo de una cadena, o su borde más largo (el mayor prefijo propio que también es sufijo)?
\item ¿Tengo en cuenta que el período mínimo puede NO dividir a la longitud, y en ese caso la cadena no es una potencia de nada salvo de sí misma?
\item ¿La cadena llega a \ten{6} caracteres y necesito un algoritmo lineal?
\end{itemize}

\graybold{Utilidad:} La función de prefijos de KMP calcula, para cada prefijo de la cadena, la longitud de su borde más largo. El borde de la cadena completa, $\pi[n-1]$, determina directamente su período mínimo $p = n - \pi[n-1]$, que es el paso más pequeño tal que $s_i = s_{i+p}$ para todo $i$ válido. Con eso se responde en tiempo lineal si una cadena es potencia de otra, cuál es su raíz primitiva, y cuántas veces se repite; es la misma herramienta que sirve para búsqueda de patrones, reutilizada sobre la propia cadena.

\graybold{Complejidad:} \bigO{n} en tiempo y memoria por cadena.

\graybold{Fundamento teórico:} Una cadena tiene período $p$ si y solo si tiene un borde de longitud $n - p$: decir que $s[0..n-p) = s[p..n)$ es exactamente decir que $s_i = s_{i+p}$ para todo $i$. Por lo tanto, el período mínimo corresponde al borde MÁS LARGO, $p = n - \pi[n-1]$. Pero tener período $p$ no basta para ser una potencia: $\texttt{abcab}$ tiene período 3 (es $\texttt{abc}$ seguido del inicio de otro $\texttt{abc}$) y sin embargo no es $a^k$ para ningún $k > 1$. Se necesita además que $p$ divida a $n$, y entonces $s = (s[0..p))^{n/p}$. Si $p$ no divide a $n$, la respuesta es 1, y no hay que buscar otro período: por el lema de periodicidad de Fine y Wilf, si $s$ tuviera un período $q$ que divide a $n$ con $q > p$, entonces $p + q \le n$ implicaría que $\gcd(p, q)$ también es período, y como $\gcd(p,q) \le p$ y $p$ es el mínimo, $\gcd(p,q) = p$, es decir, $p$ dividiría a $q$ y por lo tanto a $n$, contradicción. La función de prefijos en sí se calcula en tiempo lineal porque el borde actual $k$ crece a lo sumo en 1 por carácter y cada retroceso $k \leftarrow \pi[k-1]$ lo reduce estrictamente, de modo que el número total de retrocesos está acotado por el número total de incrementos, es decir, por $n$.~\cite{knuth1977}

\graybold{Ejercicio aplicado}

(UVA 10298 — Power Strings) Para cada cadena $s$ de la entrada, hallar el mayor $n$ tal que $s = a^n$ para alguna cadena $a$, es decir, cuántas veces se repite su bloque más corto.

\graybold{I/O:} Varias cadenas, una por línea, hasta una línea que contiene solo un punto (patrón 3 con centinela), con cadenas de hasta \ten{6} caracteres IMPRIMIBLES: eso incluye espacios, así que no se puede tokenizar con \texttt{.split()} y se recorre línea por línea con \texttt{for line in sys.stdin.buffer}, quitando solo \texttt{\textbackslash r\textbackslash n} con \texttt{rstrip} para no borrar espacios que forman parte de la cadena. Verificado contra los tres casos de ejemplo y contrastado con una fuerza bruta (probar cada divisor de la longitud) en 1508 casos, incluyendo cadenas con espacios, con signos de puntuación, la cadena \texttt{..} (que no es el centinela) y potencias truncadas cuyo período no divide la longitud, sin discrepancias. Con $|s| = 10^6$ el peor caso medido fue \aprox{}0.15s (\texttt{a...ab}, que fuerza la mayor cantidad de retrocesos), y nueve líneas de longitud máxima seguidas \aprox{}0.9s, en CPython 3.12.

\graybold{Código solución}

\begin{lstlisting}[language=Python]
import sys

def prefix_function(s):
    n = len(s)
    pi = [0]*n                  # pi[i] = borde mas largo de s[0..i]
    k = 0
    for i in range(1, n):
        c = s[i]
        # si no se puede extender el borde actual, se prueba con el
        # siguiente borde mas corto: el borde del borde
        while k and s[k] != c:
            k = pi[k-1]
        if s[k] == c:
            k += 1
        pi[i] = k
    return pi

def solve():
    out = []
    for line in sys.stdin.buffer:
        # solo se quita el salto de linea: los espacios son parte de s
        s = line.rstrip(b"\r\n")
        if s == b".":
            break
        n = len(s)
        pi = prefix_function(s)
        p = n - pi[n-1]         # periodo minimo = n - borde mas largo
        # solo es potencia si el periodo divide a la longitud (ej. "abcab" -> 1)
        out.append(str(n // p if n % p == 0 else 1))
    sys.stdout.write("\n".join(out) + "\n")

solve()
\end{lstlisting}
